\section{协作中的难点}

实现由子 agent 写，我做三件事：定目标、定验收标准、判断结论对不对。下面四类问题是我在验收里抓出来的。

\subsection{测试通过而功能没有生效}

全项目出现最多的一类问题是：代码看着对、测试全绿、功能根本没有生效。四个实例：

\begin{itemize}[leftmargin=1.5em,itemsep=0.12em,topsep=0.25em]
\item 子 agent 改服务端时漏了静态资源白名单里的新模块 \code{rules.js}，\code{/rules.js} 返回 404，\code{app.js} 的 import 图整体失败，页面元素没有渲染出来。
\item 子 agent 给陪练写了 \code{notify()}，全仓库没有调用点。主动气泡写完了，从不出现（后来补上了调用点）。
\item 子 agent 改军师时留下一个未定义变量，每秒抛一次 \code{ReferenceError}，浏览器日志里攒了 75 条。
\item 敌方补位「决定了却无人提交」：子 agent 选好了合法换宠、存进变量就返回，而玩家那条路径又拒绝替它提交，界面永久冻结在「正在出招…」。
\end{itemize}

这四次全部是子 agent 交付的缺陷，没有一次是读代码读出来的，全部是实际跑起来才看见的。

我把验收换成了两层检查。第一层是 \path{wiring.test.js}：\code{app.js} 里每个 \code{\$('id')} 都必须在 \code{index.html} 或 \code{app.js} 自建的元素里存在，死函数清单必须恰好是两个已知遗留，面试题要求的每项能力都必须在 \code{app.js} 里找得到入口；它第一次运行就找出一处指向已删除元素的残留引用。第二层是冒烟测试，用无头浏览器真打一局训练赛，断言营地渲染出伙伴卡、对局能打到结果、控制台零报错；它不进 \code{npm test}，因为需要浏览器的检查不该当成必过门。

这一层有假阳性。有一次它报「最长连续无可点 14 次」，看上去快卡死了，查下去是它从未设置播放速度，每一回合都在放动画、按钮本来就处于禁用状态。

\subsection{检查本身也要被检查}

玩家文案的术语检查写出来之后，我做了负向验证：把禁止词表里的词植入一句玩家可达的文案，测试必须失败；移除，测试必须通过。这一步让一个原本看不出来的缺口暴露了——检查取局内提示时没有传事件参数，于是「明显策略错误」那条理由对应的正文根本不在覆盖范围内。

第一次修补的方向也是错的：手工造一个事件参数传进去，负向验证立刻戳穿，把违规词放回源码，检查照样通过——那段正文要等枚举真的算出一次事故才会生成。最后换成静态扫描，逐行取 \code{coach/} 下的中文字面量对词表，不依赖任何运行时路径；它随即又扫出四处运行时检查没覆盖到的真泄漏。

\subsection{代码是对的，错的是统计}

有一处缺陷代码完全正常，错的是统计口径：玩家补位的那一回合本来就不该出招，却被计成一次「引擎兜底」，方向恰好让对手 agent 显得比实际更差。同类的一处是犹豫判定用了墙钟时间，玩家中途去操作对方面板的 13 秒被算成「犹豫了 13 秒」。这类错误不会让任何测试失败。

\subsection{哪些判断我不外包}

有一份评测结果文件存在，我打开只有两行，是更早一次局部验证的残留；此后我逐条看实际数据，不看文件名。

另一类更容易漏：子 agent 写下的「某方案没用」往往只对了一半。检索实验里最好的一臂确实没有显著赢过基线，但拆开四条臂能看到排序指标确实变好了。
